%СИСИ~Том 33   № 1-4   Год 2023





\def\stat{cont}


{%\hrule\par


%\vskip 7pt % 7pt


\raggedleft\Large \bf%\baselineskip=3.2ex


А\,В\,Т\,О\,Р\,С\,К\,И\,Й\ \ У\,К\,А\,З\,А\,Т\,Е\,Л\,Ь\ \ З\,А\ \ 2\,0\,2\,3 г. \vskip 17pt


 \hrule


 \par


\vskip 21pt plus 6pt minus 3pt }


%\vskip 11pt plus 3pt minus 2pt }





\label{st\stat}





%\def\tit{\ }





%\def\aut{\ }


%\def\auf{\ }





\def\leftkol{\ } % ENGLISH ABSTRACTS}


\def\rightkol{\ } %ENGLISH ABSTRACTS}





%\titele{\tit}{\aut}{\auf}{\leftkol}{\rightkol}





%\vspace*{-24pt}





%%%%%%%%%%%





%\thispagestyle{myheadings}


\def\leftfootline{\small{\textbf{\thepage}


\hfill СИСТЕМЫ И СРЕДСТВА ИНФОРМАТИКИ\ \ \ выпуск~33\ \ \ №~4\ \ \ 2023}


}%


 \def\rightfootline{\small{СИСТЕМЫ И СРЕДСТВА ИНФОРМАТИКИ\ \ \ выпуск~33\ \ \


№~4\ \ \ 2023


 \hfill \textbf{\thepage}}}





%%%%%%%%%%%%%


\vspace*{-1pt}





%%\vspace*{-1pt}





\noindent


{\tabcolsep=2.5pt


\begin{tabular}{p{302pt}cc}


&\multicolumn{1}{c}{\textbf{№}} & \textbf{Стр.}\\[5pt]


\Avtors{Абгарян К.\,К., Загордан Н.\,Л., Мочалова Ю.\,Д.} Метод компьютерного моделирования упругих характеристик многослойных композиционных материалов&


\raisebox{-24pt}[0pt][0pt]{4}&\raisebox{-24pt}[0pt][0pt]{\hphantom{9}92--101}\\


\Avtors{Абрамов А.\,Г., Порхачёв В.\,А.} Принципы разработки программного комплекса центра управления национальной исследовательской компьютерной сетью России&


\raisebox{-24pt}[0pt][0pt]{3}&\raisebox{-24pt}[0pt][0pt]{48--60}\\


\Avtors{Адамович~И.\,М., Волков~О.\,И.} Механизм формирования гипотез в~технологии поддержки кон\-крет\-но-ис\-то\-ри\-че\-ских исследований&


\raisebox{-24pt}[0pt][0pt]{1}&\raisebox{-24pt}[0pt][0pt]{135--145}\\


\Avtors{Адамович И.\,М., Волков О.\,И.} Несбалансированность классов в технологии поддержки кон\-крет\-но-ис\-то\-ри\-че\-ских исследований&


\raisebox{-24pt}[0pt][0pt]{4}&\raisebox{-24pt}[0pt][0pt]{149--159}\\


\Avtors{Адамович И.\,М., Волков О.\,И.} Очистка данных в~технологии поддержки кон\-крет\-но-ис\-то\-ри\-че\-ских исследований&


\raisebox{-12pt}[0pt][0pt]{3}&\raisebox{-12pt}[0pt][0pt]{149--160}\\


\Avtors{Адамович~И.\,М., Волков~О.\,И.} Применение алгоритма CHAID в~технологии поддержки конкретно-исторических исследований&


\raisebox{-24pt}[0pt][0pt]{2}&\raisebox{-24pt}[0pt][0pt]{132--141}\\


\Avtors{Архипов~П.\,О., Филиппских~С.\,Л.} Распознавание аномалий на разновременных панорамах с~использованием нейросетевого метода консолидации моделей&


\raisebox{-24pt}[0pt][0pt]{2}&\raisebox{-24pt}[0pt][0pt]{13--24}\\


\Avtors{Барашов Е.\,Б., Егоркин А.\,В., Лемтюжникова~Д.\,В., Посыпкин~М.\,А.} Анализ эффективности алгоритма редукции в~решении задачи об~упаковке в~контейнеры&


\raisebox{-24pt}[0pt][0pt]{3}&\raisebox{-24pt}[0pt][0pt]{61--75}\\


\Avtors{Бесчастный В.\,А., Голос Е.\,С., Острикова Д.\,Ю., Мачнев Е.\,А., Шоргин В.\,С., Гайдамака Ю.\,В.} Анализ совместного использования стратегий энергосбережения для устройств 5G с ограниченным функционалом&


\raisebox{-36pt}[0pt][0pt]{4}&\raisebox{-36pt}[0pt][0pt]{69--81}\\


\Avtors{Бирюкова Т.\,К., Гершкович М.\,М.} Методы оптимизации скорости выполнения функциональных запросов в автоматизированных информационных системах с учетом смыслового анализа информации&


\raisebox{-36pt}[0pt][0pt]{4}&\raisebox{-36pt}[0pt][0pt]{82--91}\\


\Avtors{Бутенко~Ю.\,И.} Использование базы данных моделей структурных переводческих трансформаций для извлечения многокомпонентных терминологических единиц&


\raisebox{-24pt}[0pt][0pt]{1}&\raisebox{-24pt}[0pt][0pt]{35--44}\\


\Avtors{Волков О.\,И.} см.\ Адамович И.\,М.&&\\


\Avtors{Волков О.\,И.} см.\ Адамович И.\,М.&&\\


\Avtors{Волков~О.\,И.} см.\ Адамович~И.\,М.&&\\


\Avtors{Волков~О.\,И.} см.\ Адамович~И.\,М.&&\\


\end{tabular}


}








\pagebreak





\def\rightkol{АВТОРСКИЙ УКАЗАТЕЛЬ ЗА 2023 г.}


\def\leftkol{АВТОРСКИЙ УКАЗАТЕЛЬ ЗА 2023 г.}





\noindent


{\tabcolsep=2.5pt


\begin{tabular}{p{302pt}cc}


&\multicolumn{1}{c}{\textbf{№}} & \textbf{Стр.}\\[6pt]


\Avtors{Гайдамака Ю.\,В.} см.\ Бесчастный В.\,А.&&\\


\Avtors{Гершкович М.\,М.} см.\ Бирюкова Т.\,К.&&\\


\Avtors{Голос Е.\,С.} см.\ Бесчастный В.\,А.&&\\


\Avtors{Гончаров А\, А.} Поиск с исключением в параллельных текстах&4&102--114\\


\Avtors{Гринченко~С.\,Н.} Базисные информационные технологии, информационная война и~многомерный иерархический территориальный суверенитет: этапы гло\-баль\-но-кос\-ми\-че\-ской коэволюции&


\raisebox{-36pt}[0pt][0pt]{1}&\raisebox{-36pt}[0pt][0pt]{146--153}\\


\Avtors{Грушо А.\,А., Грушо Н.\,А., Забежайло М.\,И., Зацаринный~А.\,А., Тимонина~Е.\,Е.} Некоторые проблемы мониторинга информационной безопасности критической инфраструктуры&


\raisebox{-36pt}[0pt][0pt]{3}&\raisebox{-36pt}[0pt][0pt]{108--116}\\


\Avtors{Грушо~А.\,А., Грушо~Н.\,А., Забежайло~М.\,И., Зацаринный~А.\,А., Тимонина~Е.\,Е., Шоргин~С.\,Я.} Задача классификации в~условиях искаженных при\-чин\-но-след\-ст\-вен\-ных связей&


\raisebox{-36pt}[0pt][0pt]{1}&\raisebox{-36pt}[0pt][0pt]{59--67}\\


\Avtors{Грушо А.\,А., Грушо Н.\,А., Забежайло М.\,И., Писковский В.\,О., Смирнов Д.\,В., Тимонина Е.\,Е.} Автоматные модели распространения сбоев и самовосстановления&


\raisebox{-24pt}[0pt][0pt]{4}&\raisebox{-24pt}[0pt][0pt]{28--37}\\


\Avtors{Грушо~А.\,А., Забежайло~М.\,И., Кульченков~В.\,В., Смирнов~Д.\,В., Тимонина~Е.\,Е., Шоргин~С.\,Я.} При\-чин\-но-след\-ст\-вен\-ные связи в~задачах анализа ненаблюдаемых свойств процессов&


\raisebox{-36pt}[0pt][0pt]{2}&\raisebox{-36pt}[0pt][0pt]{71--78}\\


\Avtors{Грушо Н.\,А.} см. Грушо А.\,А.&&\\


\Avtors{Грушо Н.\,А.} см.\ Грушо А.\,А.&&\\


\Avtors{Грушо~Н.\,А.} см.\ Грушо~А.\,А.&&\\


\Avtors{Дулин~С.\,К.} см.\ Розенберг~И.\,Н.&&\\


\Avtors{Дурново~А.\,А., Инькова~О.\,Ю., Нуриев~В.\,А.} Интеграционные возможности надкорпусных баз данных&


\raisebox{-12pt}[0pt][0pt]{1}&\raisebox{-12pt}[0pt][0pt]{24--34}\\


\Avtors{Дьяченко~Д.\,Ю.} см.\ Соколов~И.\,А.&&\\


\Avtors{Дьяченко~Д.\,Ю.} см.\ Степченков~Ю.\,А.&&\\


\Avtors{Дьяченко~Ю.\,Г.} см.\ Соколов~И.\,А.&&\\


\Avtors{Дьяченко Ю.\,Г. } см.\ Степченков Ю.\,А.&&\\


\Avtors{Дьяченко~Ю.\,Г.} см.\ Степченков~Ю.\,А.&&\\


\Avtors{Егоркин А.\,В.} см.\ Барашов Е.\,Б.&&\\


\Avtors{Егоров~В.\,Б.} Некоторые вопросы дезагрегации и~компонуемости инфраструктуры центра обработки данных&


\raisebox{-12pt}[0pt][0pt]{2}&\raisebox{-12pt}[0pt][0pt]{101--110}\\


\Avtors{Егоров~В.\,Б.} Программное определение сети в конвергентной и~гиперконвергентной инфраструктурах&


\raisebox{-12pt}[0pt][0pt]{1}&\raisebox{-12pt}[0pt][0pt]{105--113}\\


\Avtors{Егорова А.\,Ю., Зацман И.\,М., Романенко~В.\,О.} Машинный перевод с~помощью ChatGPT: мониторинг вос\-про\-из\-во\-ди\-мости результатов&


\raisebox{-24pt}[0pt][0pt]{3}&\raisebox{-24pt}[0pt][0pt]{117--128}\\


\end{tabular}


}





\pagebreak





\def\rightkol{АВТОРСКИЙ УКАЗАТЕЛЬ ЗА 2023 г.}


\def\leftkol{АВТОРСКИЙ УКАЗАТЕЛЬ ЗА 2023 г.}





\noindent


{\tabcolsep=2.5pt


\begin{tabular}{p{302pt}cc}


&\multicolumn{1}{c}{\textbf{№}} & \textbf{Стр.}\\[6pt]


\Avtors{Забежайло М.\,И.} см. Грушо А.\,А.&&\\[0.9pt]


\Avtors{Забежайло М.\,И.} см.\ Грушо А.\,А.&&\\[0.9pt]


\Avtors{Забежайло~М.\,И.} см.\ Грушо~А.\,А.&&\\[0.9pt]


\Avtors{Забежайло~М.\,И.} см.\ Грушо~А.\,А.&&\\[0.9pt]


\Avtors{Загордан Н.\,Л.} см.\ Абгарян К.\,К.&&\\[0.9pt]


\Avtors{Захаров В.\,Н.} К 40-летию выхода постановлений ЦК КПСС и~Совета Министров СССР о~дальнейшем развитии работ в~области вычислительной техники, в~том числе в~Академии наук СССР&


\raisebox{-36pt}[0pt][0pt]{3}&\raisebox{-36pt}[0pt][0pt]{\hphantom{9}4--16}\\[0.9pt]


\Avtors{Зацаринный~А.\,А., Ионенков~Ю.\,С.} Методический подход к~выбору ключевых показателей эффективности для оценки организаций ин\-фор\-ма\-ци\-он\-но-тех\-но\-ло\-ги\-че\-ской сферы&


\raisebox{-24pt}[0pt][0pt]{2}&\raisebox{-24pt}[0pt][0pt]{79--87}\\[0.9pt]


\Avtors{Зацаринный А.\,А., Сучков А.\,П.}Некоторые подходы к~анализу факторов, влияющих на~информационную безопасность систем искусственного интеллекта&


\raisebox{-24pt}[0pt][0pt]{3}&\raisebox{-24pt}[0pt][0pt]{\hphantom{9}95--107}\\[0.9pt]


\Avtors{Зацаринный А.\,А., Шабанов А.\,П.} Методический подход к~контролю качества трансформации данных в~междисциплинарной цифровой платформе&


\raisebox{-24pt}[0pt][0pt]{3}&\raisebox{-24pt}[0pt][0pt]{129--140}\\[0.9pt]


\Avtors{Зацаринный~А.\,А.} см.\ Грушо А.\,А.&&\\[0.9pt]


\Avtors{Зацаринный~А.\,А.} см.\ Грушо~А.\,А.&&\\[0.9pt]


\Avtors{Зацман И.\,М.} Научная парадигма информатики: классификация трансформаций объектов предметной области&


\raisebox{-12pt}[0pt][0pt]{4}&\raisebox{-12pt}[0pt][0pt]{126--138}\\[0.9pt]


\Avtors{Зацман И.\,М.} см.\ Егорова А.\,Ю.&&\\[0.9pt]


\Avtors{Иванов~В.\,А., Конышев~М.\,Ю., Смирнов~С.\,В., Тараканов~О.\,В., Тараканова~В.\,О., Усовик~С.\,В.} Семантические интерпретации высоких нормальных форм отношений реляционной базы данных&


\raisebox{-36pt}[0pt][0pt]{1}&\raisebox{-36pt}[0pt][0pt]{45--58}\\[0.9pt]


\Avtors{Ильин В.\,Д.} Теория S-символов: классы базовых S-задач&4&139--148\\[0.9pt]


\Avtors{Ильин~В.\,Д.} Теория S-символов: концептуальные основания&1&126--134\\[0.9pt]


\Avtors{Ильин В.\,Д.} Теория S-символов: сетевой табс-ре\-ша\-тель S-задач&3&141--148\\[0.9pt]


\Avtors{Ильин~В.\,Д.} Теория S-символов: формализация знаний об S-за\-дачах&


\raisebox{-12pt}[0pt][0pt]{2}&\raisebox{-12pt}[0pt][0pt]{124--131}\\[0.9pt]


\Avtors{Инькова~О.\,Ю.} см.\ Дурново~А.\,А.&&\\[0.9pt]


\Avtors{Ионенков~Ю.\,С.} см.\ Зацаринный~А.\,А.&&\\[0.9pt]


\Avtors{Кириков~И.\,А.} см.\ Румовская~С.\,Б.&&\\[0.9pt]


\Avtors{Ковалёв~И.\,А.} Об оценках устойчивости и~их применении для некоторых моделей массового обслуживания&


\raisebox{-12pt}[0pt][0pt]{1}&\raisebox{-12pt}[0pt][0pt]{\hphantom{1}90--104}\\[0.9pt]


\Avtors{Колин К.\,К.} Социальная эффективность информационных технологий&


\raisebox{-12pt}[0pt][0pt]{3}&\raisebox{-12pt}[0pt][0pt]{161--171}\\[0.9pt]


\Avtors{Коновалов Г.\,М.} Расчет целевых значений коэффициентов готовности для диагностик ИТЭР&


\raisebox{-12pt}[0pt][0pt]{4}&\raisebox{-12pt}[0pt][0pt]{38--49}\\


\end{tabular}


}





\pagebreak





\def\rightkol{АВТОРСКИЙ УКАЗАТЕЛЬ ЗА 2023 г.}


\def\leftkol{АВТОРСКИЙ УКАЗАТЕЛЬ ЗА 2023 г.}





\noindent


{\tabcolsep=2.5pt


\begin{tabular}{p{302pt}cc}


&\multicolumn{1}{c}{\textbf{№}} & \textbf{Стр.}\\[6pt]


\Avtors{Коновалов М.\,Г., Разумчик Р.\,В.} Алгоритм глобальной оптимизации некоторых стационарных временных характеристик заданий в~частично наблюдаемых стохастических системах с~параллельным обслуживанием&


\raisebox{-36pt}[0pt][0pt]{4}&\raisebox{-36pt}[0pt][0pt]{50--59}\\


\Avtors{Коновалов М.\,Г., Разумчик Р.\,В.}Диспетчеризация в~частично наблюдаемых стохастических системах конечной емкости с~параллельным обслуживанием&


\raisebox{-24pt}[0pt][0pt]{3}&\raisebox{-24pt}[0pt][0pt]{29--47}\\


\Avtors{Конышев~М.\,Ю.} см.\ Иванов~В.\,А.&&\\


\Avtors{Кочеткова И.\,А.} см.\ Макеева Е.\,Д.&&\\


\Avtors{Кривенко М.\,П.} Анализ монотонного тренда временного ряда&3&17--28\\


\Avtors{Кривенко~М.\,П.} Эффективные вычисления при факторизации матричных данных с~пропусками&


\raisebox{-12pt}[0pt][0pt]{1}&\raisebox{-12pt}[0pt][0pt]{78--89}\\


\Avtors{Кружков~М.\,Г.} см.\ Нуриев~В.\,А.&&\\


\Avtors{Кульченков~В.\,В.} см.\ Грушо~А.\,А.&&\\


\Avtors{Лебедев А.\,А.} см.\ Москин Н.\,Д.&&\\


\Avtors{Лемтюжникова~Д.\,В.} см.\ Барашов Е.\,Б.&&\\


\Avtors{Листопад С.\,В.} Характеристики и логическая структура методологии построения реф\-лек\-сив\-но-ак\-тив\-ных систем искусственных гетерогенных интеллектуальных агентов&


\raisebox{-24pt}[0pt][0pt]{4}&\raisebox{-24pt}[0pt][0pt]{16--27}\\


\Avtors{Макеева Е.\,Д., Кочеткова И.\,А., В.\,С. Шоргин} Модель для выбора уровней скорости широкополосного трафика eMBB в~условиях приоритетной передачи трафика URLLC в~сети~5G&


\raisebox{-36pt}[0pt][0pt]{4}&\raisebox{-36pt}[0pt][0pt]{60--68}\\


\Avtors{Мартюшова~Я.\,Г., Минеева~Т.\,А., Наумов~А.\,В.} Методы классификации пользователей СДО в~модели построения их индивидуальной траектории обучения&


\raisebox{-24pt}[0pt][0pt]{1}&\raisebox{-24pt}[0pt][0pt]{68--77}\\


\Avtors{Матвеев И.\,А.} см.\ Сафонов И.\,В.&&\\


\Avtors{Мачнев Е.\,А.} см.\ Бесчастный В.\,А.&&\\


\Avtors{Минеева~Т.\,А.} см.\ Мартюшова~Я.\,Г.&&\\


\Avtors{Морозов~Н.\,В.} см.\ Соколов~И.\,А.&&\\


\Avtors{Морозов Н.\,В. } см.\ Степченков Ю.\,А.&&\\


\Avtors{Москин Н.\,Д., Рогов А.\,А., Лебедев А.\,А.} Графовые $n$-грам\-мы в задаче атрибуции текстов&


\raisebox{-12pt}[0pt][0pt]{4}&\raisebox{-12pt}[0pt][0pt]{115--125}\\


\Avtors{Мочалова Ю.\,Д.} см.\ Абгарян К.\,К.&&\\


\Avtors{Наумов~А.\,В.} см.\ Мартюшова~Я.\,Г.&&\\


\Avtors{Никитенкова~С.\,П.} Применение биспектрального анализа в~обнаружении deepfake-изоб\-ра\-жений&


\raisebox{-12pt}[0pt][0pt]{2}&\raisebox{-12pt}[0pt][0pt]{25--33}\\


\Avtors{Никишин~Д.\,А.} Направления развития методологической базы для работы с~геоданными в~перспективных геоинформационных системах&


\raisebox{-24pt}[0pt][0pt]{2}&\raisebox{-24pt}[0pt][0pt]{34--45}\\


\Avtors{Никишин Д.\,А.} Проблемы формирования системы согласованных локальных географических онтологий&


\raisebox{-12pt}[0pt][0pt]{3}&\raisebox{-12pt}[0pt][0pt]{85--94}\\


\end{tabular}


}





\pagebreak





\def\rightkol{АВТОРСКИЙ УКАЗАТЕЛЬ ЗА 2023 г.}


\def\leftkol{АВТОРСКИЙ УКАЗАТЕЛЬ ЗА 2023 г.}





\noindent


{\tabcolsep=2.5pt


\begin{tabular}{p{302pt}cc}


&\multicolumn{1}{c}{\textbf{№}} & \textbf{Стр.}\\[6pt]\Avtors{Нуриев~В.\,А., Кружков~М.\,Г.} Корпусные данные при контрастивном изучении пунктуации&


\raisebox{-12pt}[0pt][0pt]{1}&\raisebox{-12pt}[0pt][0pt]{14--23}\\


\Avtors{Нуриев~В.\,А.} см.\ Дурново~А.\,А.&&\\


\Avtors{Орлов~Г.\,А.} см.\ Степченков~Ю.\,А.&&\\


\Avtors{Острикова Д.\,Ю.} см.\ Бесчастный В.\,А.&&\\


\Avtors{Писковский В.\,О.} см. Грушо А.\,А.&&\\


\Avtors{Плеханов Л.\,П.} см.\ Степченков Ю.\,А.&&\\


\Avtors{Порхачёв В.\,А.} см.\ Абрамов А.\,Г.&&\\


\Avtors{Посыпкин~М.\,А.} см.\ Барашов Е.\,Б.&&\\


\Avtors{Разумчик Р.\,В.} см.\ Коновалов М.\,Г.&&\\


\Avtors{Разумчик Р.\,В.} см.\ Коновалов М.\,Г.&&\\


\Avtors{Рогов А.\,А.} см.\ Москин Н.\,Д.&&\\


\Avtors{Розенберг~И.\,Н., Дулин~С.\,К.} Выбор технологических решений для поддержки процесса синтеза геоданных инфраструктуры железнодорожного транспорта&


\raisebox{-24pt}[0pt][0pt]{2}&\raisebox{-24pt}[0pt][0pt]{46--59}\\


\Avtors{Романенко~В.\,О.} см.\ Егорова А.\,Ю.&&\\


\Avtors{Румовская~С.\,Б., Кириков~И.\,А.} Визуальный язык репрезентации процесса управления конфликтами в~гибридных интеллектуальных многоагентных системах&


\raisebox{-24pt}[0pt][0pt]{2}&\raisebox{-24pt}[0pt][0pt]{60--70}\\


\Avtors{Сафонов И.\,В., Матвеев И.\,А.} Методология на\-уч\-но-ис\-сле\-до\-ва\-тель\-ской работы при~создании функций систем сканирования и~печати&


\raisebox{-24pt}[0pt][0pt]{3}&\raisebox{-24pt}[0pt][0pt]{76--84}\\


\Avtors{Смирнов Д.\,В.} см. Грушо А.\,А.&&\\


\Avtors{Смирнов~Д.\,В.} см.\ Грушо~А.\,А.&&\\


\Avtors{Смирнов~С.\,В.} см.\ Иванов~В.\,А.&&\\


\Avtors{Соколов~И.\,А., Степченков~Ю.\,А., Дьяченко~Ю.\,Г., Морозов~Н.\,В., Дьяченко~Д.\,Ю.} Самосинхронный конвейер с~переменным числом ступеней&


\raisebox{-24pt}[0pt][0pt]{1}&\raisebox{-24pt}[0pt][0pt]{\hphantom{1}4--13}\\


\Avtors{Степченков Д.\,Ю.} см.\ Степченков Ю.\,А.&&\\


\Avtors{Степченков~Д.\,Ю.} см.\ Степченков~Ю.\,А.&&\\


\Avtors{Степченков~Ю.\,А., Дьяченко~Ю.\,Г., Степченков~Д.\,Ю., Дьяченко~Д.\,Ю., Орлов~Г.\,А.} Мультиплексируемый самосинхронный конвейер&


\raisebox{-24pt}[0pt][0pt]{2}&\raisebox{-24pt}[0pt][0pt]{\hphantom{1}4--12}\\


\Avtors{Степченков Ю.\,А., Степченков Д.\,Ю., Дьяченко Ю.\,Г., Морозов Н.\,В., Плеханов Л.\,П.} Замена синхронных триггеров самосинхронными аналогами в процессе десинхронизации схемы&


\raisebox{-36pt}[0pt][0pt]{4}&\raisebox{-36pt}[0pt][0pt]{\hphantom{9}4--15}\\


\Avtors{Степченков~Ю.\,А.} см.\ Соколов~И.\,А.&&\\


\Avtors{Сучков~А.\,П.} Система ситуационного управления как мультисервисная технология в~облачной среде: сервисы аналитики&


\raisebox{-12pt}[0pt][0pt]{2}&\raisebox{-12pt}[0pt][0pt]{\hphantom{1}88--100}\\


\Avtors{Сучков~А.\,П.} Центры компетенции по искусственному интеллекту и~Национальная технологическая инициатива&


\raisebox{-12pt}[0pt][0pt]{1}&\raisebox{-12pt}[0pt][0pt]{114--125}\\


\end{tabular}


}





\pagebreak





\def\rightkol{АВТОРСКИЙ УКАЗАТЕЛЬ ЗА 2023 г.}


\def\leftkol{АВТОРСКИЙ УКАЗАТЕЛЬ ЗА 2023 г.}





\noindent


{\tabcolsep=2.5pt


\begin{tabular}{p{302pt}cc}


&\multicolumn{1}{c}{\textbf{№}} & \textbf{Стр.}\\[6pt]


\Avtors{Сучков А.\,П.} см.\ Зацаринный А.\,А.&&\\


\Avtors{Тараканов~О.\,В.} см.\ Иванов~В.\,А.&&\\


\Avtors{Тараканова~В.\,О.} см.\ Иванов~В.\,А.&&\\


\Avtors{Тимонина Е.\,Е.} см. Грушо А.\,А.&&\\


\Avtors{Тимонина~Е.\,Е.} см.\ Грушо А.\,А.&&\\


\Avtors{Тимонина~Е.\,Е.} см.\ Грушо~А.\,А.&&\\


\Avtors{Тимонина~Е.\,Е.} см.\ Грушо~А.\,А.&&\\


\Avtors{Усовик~С.\,В.} см.\ Иванов~В.\,А.&&\\


\Avtors{Филиппских~С.\,Л.} см.\ Архипов~П.\,О&&\\


\Avtors{Читалов~Д.\,И.} Разработка модуля для работы с~решателем plasticStressedFoam на базе пакета OpenFOAM&


\raisebox{-12pt}[0pt][0pt]{2}&\raisebox{-12pt}[0pt][0pt]{111--123}\\


\Avtors{Шабанов А.\,П.} см.\ Зацаринный А.\,А.&&\\


\Avtors{Шоргин В.\,С.} см.\ Бесчастный В.\,А.&&\\


\Avtors{Шоргин В.\,С.} см.\ Макеева Е.\,Д.&&\\


\Avtors{Шоргин~С.\,Я.} см.\ Грушо~А.\,А.&&\\


\Avtors{Шоргин~С.\,Я.} см.\ Грушо~А.\,А.&&\\


\end{tabular}


}








%\thispagestyle{myheadings}


\def\leftfootline{\small{\textbf{\thepage}


\hfill СИСТЕМЫ И СРЕДСТВА ИНФОРМАТИКИ\ \ \ выпуск~33\ \ \ №~4\ \ \ 2023}


}%


 \def\rightfootline{\small{СИСТЕМЫ И СРЕДСТВА ИНФОРМАТИКИ\ \ \ выпуск~33\ \ \


№~4\ \ \ 2023


 \hfill \textbf{\thepage}}}


 \label{end\stat}





\newpage





\def\stat{cont-e}


{%\hrule\par


%\vskip 7pt % 7pt


\raggedleft\Large \bf%\baselineskip=3.2ex


2\,0\,2\,3\ \ A\,U\,T\,H\,O\,R\ \ I\,N\,D\,E\,X \vskip 17pt


 \hrule


 \par


%\vskip 21pt plus 6pt minus 3pt }


\vskip 11pt plus 3pt minus 2pt }





\label{st\stat}





%\def\tit{\ }





%\def\aut{\ }


%\def\auf{\ }





\def\leftkol{\ } %2010 AUTHOR INDEX} % ENGLISH ABSTRACTS}


\def\rightkol{\ } %2010 AUTHOR INDEX} %ENGLISH ABSTRACTS}





%\titele{\tit}{\aut}{\auf}{\leftkol}{\rightkol}





%\vspace*{-12pt}





%%%%%%%%%%%





\def\leftfootline{\small{\textbf{\thepage}


\hfill Sistemy i Sredstva Informatiki~---


Systems and Means of Informatics 2023 vol~33 no\,4}


}%


 \def\rightfootline{\small{Sistemy i Sredstva Informatiki~---


 Systems and Means of Informatics 2023 vol~33 no\,4


\hfill \textbf{\thepage}}}





\noindent


{\tabcolsep=3pt


\begin{tabular}{p{296pt}cc}


&\multicolumn{1}{c}{\textbf{No.}} & \textbf{Page}\\[5pt]


\Avtors{Abgaryan K.\,K., Zagordan N.\,L., and Mochalova Y.\,D.} Method for computer modeling of elastic characteristics of multilayer composite materials&


\raisebox{-24pt}[0pt][0pt]{4}&\raisebox{-24pt}[0pt][0pt]{\hphantom{9}92--101}\\


\Avtors{Abramov A.\,G. and Porkhachev V.\,A.} Principles of development of~a~software suite for~the~network operations Center of~the~National Research Computer Network of~Russia&


\raisebox{-24pt}[0pt][0pt]{3}&\raisebox{-24pt}[0pt][0pt]{48--60}\\


\Avtors{Adamovich~I.\,M. and~Volkov~O.\,I.} Application of the CHAID algorithm in the technology of concrete historical investigation support&


\raisebox{-24pt}[0pt][0pt]{2}&\raisebox{-24pt}[0pt][0pt]{132--141}\\


\Avtors{Adamovich I.\,M. and Volkov O.\,I.} Class imbalance in the technology of concrete historical investigation support&


\raisebox{-12pt}[0pt][0pt]{4}&\raisebox{-12pt}[0pt][0pt]{149--159}\\


\Avtors{Adamovich I.\,M. and Volkov O.\,I.} Data cleansing in~the~technology of~concrete historical investigation support&


\raisebox{-12pt}[0pt][0pt]{3}&\raisebox{-12pt}[0pt][0pt]{149--160}\\


\Avtors{Adamovich~I.\,M. and~Volkov~O.\,I.} Hypothesis formation mechanism in the technology of concrete historical investigation support&


\raisebox{-24pt}[0pt][0pt]{1}&\raisebox{-24pt}[0pt][0pt]{135--145}\\


\Avtors{Arkhipov~P.\,O. and~Philippskih~S.\,L.} Recognition of anomalies on multitime panoramas using the neural network method of model amalgamation&


\raisebox{-24pt}[0pt][0pt]{2}&\raisebox{-24pt}[0pt][0pt]{13--24}\\


\Avtors{Barashov E.\,B., Egorkin A.\,V., Lemtyuzhnikova D.\,V., and~Posypkin~M.\,A.} Efficiency of~the~reduction algorithms in~the~bin packing problem&


\raisebox{-24pt}[0pt][0pt]{3}&\raisebox{-24pt}[0pt][0pt]{61--75}\\


\Avtors{Beschastnyi V.\,A., Golos E.\,S., Ostrikova D.\,Yu., Machnev~E.\,A., Shorgin V.\,S., and Gaidamaka Yu.\,V.} Analysis of joint usage of energy conservation strategies for 5G devices with reduced capability&


\raisebox{-36pt}[0pt][0pt]{4}&\raisebox{-36pt}[0pt][0pt]{69--81}\\


\Avtors{Biryukova T.\,K. and Gershkovich M.\,M.} Methods for optimizing the speed of performing functional requests in automated information systems taking into account the semantic analysis of information&


\raisebox{-36pt}[0pt][0pt]{4}&\raisebox{-36pt}[0pt][0pt]{82--91}\\


\Avtors{Butenko~Yu.\,I.} Using a~database of structural transformations to extract multicomponent terminological units&


\raisebox{-12pt}[0pt][0pt]{1}&\raisebox{-12pt}[0pt][0pt]{35--44}\\


\Avtors{Chitalov~D.\,I.} Development of a~module for working with the plasticStressedFoam solver based on the OpenFOAM package&


\raisebox{-12pt}[0pt][0pt]{2}&\raisebox{-12pt}[0pt][0pt]{111--123}\\


\Avtors{Diachenko~D.\,Yu.} see Sokolov~I.\,A.&&\\


\Avtors{Diachenko~D.\,Yu.} see Stepchenkov~Yu.\,A.&&\\


\Avtors{Diachenko~Yu.\,G.} see Sokolov~I.\,A.&&\\


\Avtors{Diachenko Yu.\,G.} see Stepchenkov Yu.\,A.&&\\


\end{tabular}


}





\pagebreak





\def\rightkol{2023 AUTHOR INDEX}


\def\leftkol{2023 AUTHOR INDEX}





\noindent


{\tabcolsep=3pt


\begin{tabular}{p{296pt}cc}


&\multicolumn{1}{c}{\textbf{No.}} & \textbf{Page}\\[6pt]


\Avtors{Diachenko~Yu.\,G.} see Stepchenkov~Yu.\,A.&&\\


\Avtors{Dulin~S.\,K.} see Rozenberg~I.\,N.&&\\


\Avtors{Durnovo~A.\,A., Inkova~O.\,Yu., and~Nuriev~V.\,A.} Integration capacities of supracorpora databases&


\raisebox{-12pt}[0pt][0pt]{1}&\raisebox{-12pt}[0pt][0pt]{24--34}\\


\Avtors{Egorkin A.\,V.} see Barashov E.\,B.&&\\


\Avtors{Egorov~V.\,B.} Some issues of disaggregation and composability of the data center infrastructure&


\raisebox{-12pt}[0pt][0pt]{2}&\raisebox{-12pt}[0pt][0pt]{101--110}\\


\Avtors{Egorov~V.\,B.} The software-defined networking in converged and hyperconverged infrastructures&


\raisebox{-12pt}[0pt][0pt]{1}&\raisebox{-12pt}[0pt][0pt]{105--113}\\


\Avtors{Egorova A.\,Yu., Zatsman I.\,M., and~Romanenko~V.\,O.} Machine translation by~ChatGPT: Monitoring of~outcome reproducibility&


\raisebox{-24pt}[0pt][0pt]{3}&\raisebox{-24pt}[0pt][0pt]{117--128}\\


\Avtors{Gaidamaka Yu.\,V.} see Beschastnyi V.\,A.&&\\


\Avtors{Gershkovich M.\,M.} see Biryukova T.\,K.&&\\


\Avtors{Golos E.\,S.} see Beschastnyi V.\,A.&&\\


\Avtors{Goncharov A.\,A.} Search with exclusion in parallel texts&4&102--114\\


\Avtors{Grinchenko~S.\,N.} Basic information technologies, information warfare, and multidimensional hierarchical territorial sovereignty: Stages of global-space coevolution&


\raisebox{-24pt}[0pt][0pt]{1}&\raisebox{-24pt}[0pt][0pt]{146--153}\\


\Avtors{Grusho A.\,A., Grusho N.\,A., Zabezhailo M.\,I., Pis\-kov\-ski~V.\,O., Smirnov D.\,V., and Timonina E.\,E.} Automata models of fault propagation and self-healing&


\raisebox{-24pt}[0pt][0pt]{4}&\raisebox{-24pt}[0pt][0pt]{28--37}\\


\Avtors{Grusho A.\,A., Grusho N.\,A., Zabezhailo M.\,I., Zatsarinny~A.\,A., and~Timonina~E.\,E.} Some challenges of~critical infrastructure information security monitoring&


\raisebox{-24pt}[0pt][0pt]{3}&\raisebox{-24pt}[0pt][0pt]{108--116}\\


\Avtors{Grusho~A.\,A., Grusho~N.\,A., Zabezhailo~M.\,I., Zatsarinny~A.\,A., Timonina~E.\,E., and~Shorgin~S.\,Ya.} Classification problem in conditions of distorted cause-and-effect relationships&


\raisebox{-36pt}[0pt][0pt]{1}&\raisebox{-36pt}[0pt][0pt]{59--67}\\


\Avtors{Grusho~A.\,A., Zabezhailo~M.\,I., Kulchenkov~V.\,V., Smirnov~D.\,V., Timonina~E.\,E., and~Shorgin~S.\,Ya.} Cause-and-effect relationships in analysis of unobservable process properties&


\raisebox{-36pt}[0pt][0pt]{2}&\raisebox{-36pt}[0pt][0pt]{71--78}\\


\Avtors{Grusho~N.\,A.} see Grusho~A.\,A.&&\\


\Avtors{Grusho N.\,A.} see Grusho A.\,A.&&\\


\Avtors{Grusho N.\,A.} see Grusho A.\,A.&&\\


\Avtors{Ilyin V.\,D.} Theory of S-symbols: Network tabs-solver of \mbox{S-problems}&


\raisebox{-12pt}[0pt][0pt]{3}&\raisebox{-12pt}[0pt][0pt]{141--148}\\


\Avtors{Ilyin~V.\,D.} Theory of S-symbols: Conceptual foundations&1&126--134\\


\Avtors{Ilyin~V.\,D.} Theory of S-symbols: Formalization of knowledge about S-problems&


\raisebox{-12pt}[0pt][0pt]{2}&\raisebox{-12pt}[0pt][0pt]{124--131}\\


\Avtors{Ilyin V.\,D.} Theory of S-symbols: The classes of basic S-problems&4&139--148\\


\end{tabular}


}





\pagebreak





\def\rightkol{2023 AUTHOR INDEX}


\def\leftkol{2023 AUTHOR INDEX}





\noindent


{\tabcolsep=3pt


\begin{tabular}{p{296pt}cc}


&\multicolumn{1}{c}{\textbf{No.}} & \textbf{Page}\\[6pt]


\Avtors{Inkova~O.\,Yu.} see Durnovo~A.\,A.&&\\


\Avtors{Ionenkov~Yu.\,S.} see Zatsarinny~A.\,A.&&\\


\Avtors{Ivanov~V.\,A., Konyshev~M.\,Yu., Smirnov~S.\,V., Tarakanov~O.\,V., Tarakanova~V.\,O., and~Usovik~S.\,V.} Semantic interpretations of high normal forms of relations in a~relational database&


\raisebox{-36pt}[0pt][0pt]{1}&\raisebox{-36pt}[0pt][0pt]{45--58}\\


\Avtors{Kirikov~I.\,A.} see Rumovskaya~S.\,B.&&\\


\Avtors{Kochetkova I.\,A.} see Makeeva E.\,D.&&\\


\Avtors{Kolin K.\,K.} Social efficiency of~information technologies&3&161--171\\


\Avtors{Konovalov G.\,M.} Calculation of availability targets for ITER diagnostics&


\raisebox{-12pt}[0pt][0pt]{4}&\raisebox{-12pt}[0pt][0pt]{38--49}\\


\Avtors{Konovalov M.\,G. and Razumchik R.\,V.} Algorithm for global optimization of time-related stationary characteristics of jobs in nonobservable parallel queues&


\raisebox{-24pt}[0pt][0pt]{4}&\raisebox{-24pt}[0pt][0pt]{50--59}\\


\Avtors{Konovalov M.\,G. and~Razumchik R.\,V.} Dispatching in~nonobservable parallel queues with~finite capacities&


\raisebox{-12pt}[0pt][0pt]{3}&\raisebox{-12pt}[0pt][0pt]{29--47}\\


\Avtors{Konyshev~M.\,Yu.} see Ivanov~V.\,A.&&\\


\Avtors{Kovalev~I.\,A.} On the perturbation bounds and their application for some queueing models&


\raisebox{-12pt}[0pt][0pt]{1}&\raisebox{-12pt}[0pt][0pt]{\hphantom{1}90--104}\\


\Avtors{Krivenko~M.\,P.} Efficient computations in matrix factorization with missing components&


\raisebox{-12pt}[0pt][0pt]{1}&\raisebox{-12pt}[0pt][0pt]{78--89}\\


\Avtors{Krivenko M.\,P.} Time series monotonic trend analysis&3&17--28\\


\Avtors{Kruzhkov~M.\,G.} see Nuriev~V.\,A.&&\\


\Avtors{Kulchenkov~V.\,V.} see Grusho~A.\,A.&&\\


\Avtors{Lebedev A.\,A.} see Moskin N.\,D.&&\\


\Avtors{Lemtyuzhnikova D.\,V.} see Barashov E.\,B.&&\\


\Avtors{Listopad S.\,V.} Characteristics and logical structure of the methodology for constructing reflexive-active systems of artificial heterogeneous intelligent agents&


\raisebox{-24pt}[0pt][0pt]{4}&\raisebox{-24pt}[0pt][0pt]{16--27}\\


\Avtors{Machnev E.\,A.} see Beschastnyi V.\,A.&&\\


\Avtors{Makeeva E.\,D., Kochetkova I.\,A., and Shorgin V.\,S.} Queuing model for choosing eMBB bit rate levels under URLLC priority transmission in a 5G network&


\raisebox{-24pt}[0pt][0pt]{4}&\raisebox{-24pt}[0pt][0pt]{60--68}\\


\Avtors{Martyushova~Ya.\,G., Mineyeva~T.\,A., and~Naumov~A.\,V.} Methods of classifying the distance learning system users in the model of constructing their personalized learning strategies&


\raisebox{-36pt}[0pt][0pt]{1}&\raisebox{-36pt}[0pt][0pt]{68--77}\\


\Avtors{Matveev I.\,A.} see Safonov I.\,V.&&\\


\Avtors{Mineyeva~T.\,A.} see Martyushova~Ya.\,G.&&\\


\Avtors{Mochalova Y.\,D.} see Abgaryan K.\,K.&&\\


\Avtors{Morozov~N.\,V.} see Sokolov~I.\,A.&&\\


\Avtors{Morozov N.\,V.} see Stepchenkov Yu.\,A.&&\\


\end{tabular}


}





\pagebreak





\def\rightkol{2023 AUTHOR INDEX}


\def\leftkol{2023 AUTHOR INDEX}





\noindent


{\tabcolsep=3pt


\begin{tabular}{p{296pt}cc}


&\multicolumn{1}{c}{\textbf{No.}} & \textbf{Page}\\[6pt]


\Avtors{Moskin N.\,D., Rogov A.\,A., and Lebedev A.\,A.} Graph \mbox{$n$-grams} in the text attribution problem&


\raisebox{-12pt}[0pt][0pt]{4}&\raisebox{-12pt}[0pt][0pt]{115--125}\\[1pt]


\Avtors{Naumov~A.\,V.} see Martyushova~Ya.\,G.&&\\[1pt]


\Avtors{Nikishin~D.\,A.} Directions of development of the methodological base for working with geodata in promising geoinformation systems&


\raisebox{-24pt}[0pt][0pt]{2}&\raisebox{-24pt}[0pt][0pt]{34--45}\\[1pt]


\Avtors{Nikishin D.\,A.} Problems of formation of~a~system of coordinated local geographical ontologies&


\raisebox{-12pt}[0pt][0pt]{3}&\raisebox{-12pt}[0pt][0pt]{85--94}\\[1pt]


\Avtors{Nikitenkova~S.\,P.} Deepfake image detection using bispectral analysis&


\raisebox{-12pt}[0pt][0pt]{2}&\raisebox{-12pt}[0pt][0pt]{25--33}\\[1pt]


\Avtors{Nuriev~V.\,A. and~Kruzhkov~M.\,G.} The parallel corpora perspective on studying contrastive punctuation&


\raisebox{-12pt}[0pt][0pt]{1}&\raisebox{-12pt}[0pt][0pt]{14--23}\\[1pt]


\Avtors{Nuriev~V.\,A.} see Durnovo~A.\,A.&&\\[1pt]


\Avtors{Orlov~G.\,A.} see Stepchenkov~Yu.\,A.&&\\[1pt]


\Avtors{Ostrikova D.\,Yu.} see Beschastnyi V.\,A.&&\\[1pt]


\Avtors{Philippskih~S.\,L.} see Arkhipov~P.\,O.&&\\[1pt]


\Avtors{Piskovski V.\,O.} see Grusho A.\,A.&&\\[1pt]


\Avtors{Plekhanov L.\,P.} see Stepchenkov Yu.\,A.&&\\[1pt]


\Avtors{Porkhachev V.\,A.} see Abramov A.\,G.&&\\[1pt]


\Avtors{Posypkin~M.\,A.} see Barashov E.\,B.&&\\[1pt]


\Avtors{Razumchik R.\,V.} see Konovalov M.\,G.&&\\[1pt]


\Avtors{Razumchik R.\,V.} see Konovalov M.\,G.&&\\[1pt]


\Avtors{Rogov A.\,A.} see Moskin N.\,D.&&\\[1pt]


\Avtors{Romanenko~V.\,O.} see Egorova A.\,Yu.&&\\[1pt]


\Avtors{Rozenberg~I.\,N. and~Dulin~S.\,K.} Selection of technological solutions to support the process of synthesis of geodata of railway transport infrastructure&


\raisebox{-24pt}[0pt][0pt]{2}&\raisebox{-24pt}[0pt][0pt]{46--59}\\[1pt]


\Avtors{Rumovskaya~S.\,B. and~Kirikov~I.\,A.} A visual language of the representation of the conflict management process in hybrid intelligent multiagent systems&


\raisebox{-24pt}[0pt][0pt]{2}&\raisebox{-24pt}[0pt][0pt]{60--70}\\[1pt]


\Avtors{Safonov I.\,V. and Matveev I.\,A.} Research methodology for~creating functions of~scanning and~printing systems&


\raisebox{-12pt}[0pt][0pt]{3}&\raisebox{-12pt}[0pt][0pt]{76--84}\\[1pt]


\Avtors{Shabanov A.\,P.} see Zatsarinny A.\,A.&&\\[1pt]


\Avtors{Shorgin~S.\,Ya.} see Grusho~A.\,A.&&\\[1pt]


\Avtors{Shorgin~S.\,Ya.} see Grusho~A.\,A.&&\\[1pt]


\Avtors{Shorgin V.\,S.} see Beschastnyi V.\,A.&&\\[1pt]


\Avtors{Shorgin V.\,S.} see Makeeva E.\,D.&&\\[1pt]


\Avtors{Smirnov D.\,V.} see Grusho A.\,A.&&\\[1pt]


\Avtors{Smirnov~D.\,V.} see Grusho~A.\,A.&&\\[1pt]


\Avtors{Smirnov~S.\,V.} see Ivanov~V.\,A.&&\\


\end{tabular}


}





\pagebreak





\def\rightkol{2023 AUTHOR INDEX}


\def\leftkol{2023 AUTHOR INDEX}





\noindent


{\tabcolsep=3pt


\begin{tabular}{p{296pt}cc}


&\multicolumn{1}{c}{\textbf{No.}} & \textbf{Page}\\[6pt]


\Avtors{Sokolov~I.\,A., Stepchenkov~Yu.\,A., Diachenko~Yu.\,G., Morozov~N.\,V., and~Diachenko~D.\,Yu.} Self-timed pipeline with variable stage number&


\raisebox{-24pt}[0pt][0pt]{1}&\raisebox{-24pt}[0pt][0pt]{\hphantom{1}4--13}\\


\Avtors{Stepchenkov D.\,Yu.} see Stepchenkov Yu.\,A.&&\\


\Avtors{Stepchenkov~D.\,Yu.} see Stepchenkov~Yu.\,A.&&\\


\Avtors{Stepchenkov~Yu.\,A., Diachenko~Yu.\,G., Stepchenkov~D.\,Yu., Diachenko~D.\,Yu., and~Orlov~G.\,A.} Multiplexed self-timed pipeline&


\raisebox{-24pt}[0pt][0pt]{2}&\raisebox{-24pt}[0pt][0pt]{\hphantom{1}4--12}\\


\Avtors{Stepchenkov Yu.\,A., Stepchenkov D.\,Yu., Diachenko Yu.\,G., Morozov N.\,V., and Plekhanov L.\,P.} Replacing synchronous triggers with self-timed counterparts during circuit desynchronization&


\raisebox{-36pt}[0pt][0pt]{4}&\raisebox{-36pt}[0pt][0pt]{\hphantom{9}4--15}\\


\Avtors{Stepchenkov~Yu.\,A.} see Sokolov~I.\,A.&&\\


\Avtors{Suchkov~A.\,P.} Competence centers for artificial intelligence and the National Technology Initiative&


\raisebox{-12pt}[0pt][0pt]{1}&\raisebox{-12pt}[0pt][0pt]{114--125}\\


\Avtors{Suchkov~A.\,P.} Situational management system as a~multiservice technology in a~cloud environment: Analytic services&


\raisebox{-12pt}[0pt][0pt]{2}&\raisebox{-12pt}[0pt][0pt]{\hphantom{1}88--100}\\


\Avtors{Suchkov A.\,P.} see Zatsarinny A.\,A.&&\\


\Avtors{Tarakanov~O.\,V.} see Ivanov~V.\,A.&&\\


\Avtors{Tarakanova~V.\,O.} see Ivanov~V.\,A.&&\\


\Avtors{Timonina E.\,E.} see Grusho A.\,A.&&\\


\Avtors{Timonina~E.\,E.} see Grusho A.\,A.&&\\


\Avtors{Timonina~E.\,E.} see Grusho~A.\,A.&&\\


\Avtors{Timonina~E.\,E.} see Grusho~A.\,A.&&\\


\Avtors{Usovik~S.\,V.} see Ivanov~V.\,A.&&\\


\Avtors{Volkov O.\,I.} see Adamovich I.\,M.&&\\


\Avtors{Volkov O.\,I.} see Adamovich I.\,M.&&\\


\Avtors{Volkov~O.\,I.} see Adamovich~I.\,M.&&\\


\Avtors{Volkov~O.\,I.} see Adamovich~I.\,M.&&\\


\Avtors{Zabezhailo M.\,I.} see Grusho A.\,A. &&\\


\Avtors{Zabezhailo M.\,I.} see Grusho A.\,A.&&\\


\Avtors{Zabezhailo~M.\,I.} see Grusho~A.\,A.&&\\


\Avtors{Zabezhailo~M.\,I.} see Grusho~A.\,A.&&\\


\Avtors{Zagordan N.\,L.} see Abgaryan K.\,K.&&\\


\Avtors{Zakharov V.\,N.} To the 40th Anniversary of~the~Decrees of~the~Central Committee of~the~CPSU and~the~Council of~Ministers of~the~USSR on~the~further development of~work in~the~field of~computer technology, including at~the~USSR Academy of Sciences&


\raisebox{-48pt}[0pt][0pt]{3}&\raisebox{-48pt}[0pt][0pt]{\hphantom{9}4--16}\\


\Avtors{Zatsarinny~A.\,A. and~Ionenkov~Yu.\,S.} Methodological approach to the selection of key performance indicators for evaluating the information and technology organizations&


\raisebox{-24pt}[0pt][0pt]{2}&\raisebox{-24pt}[0pt][0pt]{79--87}\\


\end{tabular}


}





\pagebreak





\def\rightkol{2023 AUTHOR INDEX}


\def\leftkol{2023 AUTHOR INDEX}





\noindent


{\tabcolsep=3pt


\begin{tabular}{p{296pt}cc}


&\multicolumn{1}{c}{\textbf{No.}} & \textbf{Page}\\[6pt]


\Avtors{Zatsarinny A.\,A. and Shabanov A.\,P.} Methodological approach to~quality control of~data transformation in~an~interdisciplinary digital platform&


\raisebox{-24pt}[0pt][0pt]{3}&\raisebox{-24pt}[0pt][0pt]{129--140}\\


\Avtors{Zatsarinny A.\,A. and Suchkov A.\,P.} Some approaches to~the~analysis of factors affecting the~information security of~artificial intelligence systems&


\raisebox{-24pt}[0pt][0pt]{3}&\raisebox{-24pt}[0pt][0pt]{\hphantom{9}95--107}\\


\Avtors{Zatsarinny~A.\,A.} see Grusho A.\,A.&&\\


\Avtors{Zatsarinny~A.\,A.} see Grusho~A.\,A.&&\\


\Avtors{Zatsman I.\,M.} Scientific paradigm of informatics: Transformation classification of domain objects&


\raisebox{-12pt}[0pt][0pt]{4}&\raisebox{-12pt}[0pt][0pt]{126--138}\\


\Avtors{Zatsman I.\,M.} see Egorova A.\,Yu.&&\\


\end{tabular}


}





\pagebreak





\def\rightkol{2023 AUTHOR INDEX}


\def\leftkol{2023 AUTHOR INDEX}








\def\leftfootline{\small{\textbf{\thepage}


\hfill Sistemy i Sredstva Informatiki~---


Systems and Means of Informatics 2023 vol~33 no\,4}


}%


 \def\rightfootline{\small{Sistemy i Sredstva Informatiki~---


 Systems and Means of Informatics 2023 vol~33 no\,4


\hfill \textbf{\thepage}}}


 \label{end\stat}


